% !TEX TS-program = xelatex --shell-escape
% !TEX encoding = UTF-8 Unicode
%
\documentclass{beamer}
%\setbeamertemplate{items}[ball] 
\DeclareMathOperator*{\argmax}{arg\,max}
\usecolortheme{whale}
\useinnertheme[shadow]{rounded}
\usenavigationsymbolstemplate{}
\usepackage[subsection=false,headline=empty,footline=outlineauthortitle]{beamerouterthememiniframesbottom}
%\usepackage{beamerthemesplit}
\usepackage{tikz}
\usepackage{pgfplots}
\usetikzlibrary{arrows,automata,positioning,shapes,shapes.geometric,fit,tikzmark}
\usefonttheme[onlymath]{serif}

\usepackage{enumitem}

\usepackage{textcomp}
% Set up citation style
\usepackage{natbib}
\usepackage{bibentry}
\bibpunct{(}{)}{;}{a}{,}{,}
\newcommand{\newcite}[1]{\citet{#1}}
\renewcommand{\cite}[1]{\citep{#1}}

%\setbeamercovered{transparent}

% Specify that no bibliography should be printed
\bibliographystyle{plainnat}


\setbeamertemplate{footline}[frame number]{}

%gets rid of bottom navigation symbols
\setbeamertemplate{navigation symbols}{}

%gets rid of footer
%will override 'frame number' instruction above
%comment out to revert to previous/default definitions
\setbeamertemplate{footline}{}

\usepackage{pgfpages}
\usepackage{graphicx}
\usepackage{enumitem}
\setlist[itemize,1]{label={\fontfamily{cmr}\fontencoding{T1}\selectfont\textbullet}}


\usepackage{pgfplots}

%\usepackage{ dsfont }
%\setbeamertemplate{note page}[plain]
%\setbeameroption{show notes on second screen=right}
%\usepackage{minted}
\title{Neural Network Concepts}
%\author{Lane Schwartz}

%\institute[shortinst]{University of Illinois at Urbana-Champaign}

\author{Differential calculus}

\date{}

\begin{document}

\frame{
\titlepage

}

\frame{

\begin{center}

Chain rule

\ \\

\begin{align*}
x &= \ldots \\
y &= f(x) \\
z &= g(y)
\end{align*}

\end{center}

}


\frame{

\begin{center}

Let $f$ and $g$ each be functions with signature $\mathrm{I\!R} \rightarrow \mathrm{I\!R}$

\pause
\ \\

Now, define:

\begin{equation}
h(x) = f(x) + g(x)
\end{equation}

\end{center}

}


\frame{

Definition of derivative:

\begin{align*}
f^{\prime}(x) \defeq \Lim{\Delta x \rightarrow 0} \frac{f(x + \Delta x) - f(x) }{\Delta x}
\end{align*}

}

\frame{ 

\begin{align*}
h(x) = f(x) + g(x)
\end{align*}

\begin{align*}
\onslide<2->{h^{\prime}(x) &= \Lim{\Delta x \rightarrow 0} \frac{h(x + \Delta x) - h(x) }{\Delta x} \\ }
\onslide<3->{ &= \Lim{\Delta x \rightarrow 0} \frac{\left[ f(x + \Delta x) + g(x + \Delta x) \right]  - [f(x) + g(x)] }{\Delta x} \\ }
\onslide<4->{ &= \Lim{\Delta x \rightarrow 0} \frac{\left[ f(x + \Delta x) + g(x + \Delta x) \right]  - f(x) - g(x)] }{\Delta x}  \\ }
\onslide<5->{ &= \Lim{\Delta x \rightarrow 0} \frac{[ f(x + \Delta x) - f(x)] + [ g(x + \Delta x) - g(x) ] }{\Delta x}  \\ }
\onslide<6->{ &= \Lim{\Delta x \rightarrow 0} \frac{f(x + \Delta x) - f(x)}{\Delta x} + \Lim{\Delta x \rightarrow 0} \frac{g(x + \Delta x) - g(x)}{\Delta x}  \\ }
\onslide<7->{h^{\prime}(x) &= f^{\prime}(x) + g^{\prime}(x) \\ }
\end{align*}

}

\frame{ 

\begin{align*}
h(x) &= f(x) + g(x) \\
h^{\prime}(x) &= f^{\prime}(x) + g^{\prime}(x)
\end{align*}

\pause
Now, let $c$ be a constant, and define $f$ and $g$ as:
\pause
\begin{align*}
f(x) &= x \ & g(x) &= c\\
f^{\prime}(x) &= 1   & g^{\prime}(x) &= 0 \ 
\end{align*}

\pause
Therefore

\begin{align*}
h(x) &= f(x) + g(x) \\
     &= x + c \\
h^{\prime}(x) &= f^{\prime}(x) + g^{\prime}(x) \\
              &= 1 + 0 \\
 &= 1 \\
\end{align*}

}

\frame{

\begin{align*}
f(x) &= x + c \\
f^{\prime}(x) &= 1 \\
\end{align*}

}

\frame{

Definition of partial derivative:

\begin{align*}
f_{x_i}^{\prime}(x_1, \ldots, x_{i-1}, x_i, x_{i+1} \ldots, x_n) \\ \ \\
\frac{\partial}{\partial x_i} f(x_1, \ldots, x_{i-1}, x_i, x_{i+1} \ldots, x_n)
%\defeq \Lim{\Delta x \rightarrow 0} \frac{f(x + \Delta x) - f(x) }{\Delta x}
\end{align*}

\begin{align*}
f_{x_i} &= \frac{\partial f}{\partial x_i} \defeq \\ \ \\
& \Lim{\Delta x_i \rightarrow 0} \frac{f(x_1, \ldots, x_{i-1}, x_i  + \Delta x_i, x_{i+1} \ldots, x_n) - f(x_1, \ldots, x_n) }{\Delta x_i}
\end{align*}

}


\frame{

\begin{align*}
f(x_1, x_2) &= x_1 + x_2 \\
f_{x_1}^{\prime}(x_1, x_2) &= 1 \\
\frac{\partial}{\partial x} f(x_1, x_2) &= 1 \\
\end{align*}

}

\frame{


  \begin{center}
  \begin{tikzpicture}
    [neuron/.style={circle,draw=black,fill=white,inner sep=0pt,minimum size=10mm,outer sep=0.5pt},
    >=stealth',
    scale=0.7,transform shape]
    
    \node          (i01)                              {$x_1$};
    \node          (b00)      [above=10mm of i01]     {$1$};
    
    \node [neuron] (n01)      [right=20mm of i01]     {$\tanh$};
%    \node [neuron] (b10)      [above=10mm of n01]     {$1$};
%    \node          (inv)      [below=5mm of  n01]     {\ };
%    \node [neuron] (n02)      [below=10mm of n01]     {$f_4$};

    \node          (i02)      [below=10mm of i01]      {$x_2$};
    

%    \node [neuron] (n03)      [right=20mm of n01]     {$f_3$};
%    \node          (inv)      [below=5mm of  n03]     {\ };
%    \node [neuron] (n04)      [right=20mm of n02]     {$f_4$};

%    \node [neuron] (n03)      [right=20mm of inv]     {$f_{14}$};
    
%    \node (out)      [right=20mm of n01] {$o$};
    
    \draw [->] (b00) -- (n01) node[sloped, midway, above] {$w_3$} ;
    \draw [->] (i01) -- (n01) node[midway, above] {$w_1$} ;
    \draw [->] (i02) -- (n01) node[sloped, midway, above] {$w_2$} ;
    
%    \draw [->] (b00) -- (n02);
%    \draw [->] (i01) -- (n02);
%    \draw [->] (i02) -- (n02);
%    
%    \draw [->] (n01) -- (n03);
%    \draw [->] (n02) -- (n03);
%    
%    \draw [->] (n01) -- (n04);
%    \draw [->] (n02) -- (n04);
%    
%    \draw [->] (n01) -- (n03);
%    \draw [->] (n02) -- (n03);
    
%    \draw [->] (n01) -- (out);
    
  \end{tikzpicture}

  \end{center}



\begin{align*}
f(x_1, x_2) &= \tanh(x_1 \times w_1 + x_2 \times w_2) \\
\end{align*}

}




\begin{frame}

\note{
This presentation was created and narrated by Lane Schwartz. \\ \ \\
You are free to reproduce and adapt this work under the terms of the Creative Commons Attribution-ShareAlike 4.0 International License.
}

\begin{center}
\includegraphics[scale=0.3]{images/CreativeCommons} \ \ \ \ \ 
\includegraphics[scale=0.3]{images/CreativeCommonsAttribution} \ \ \ \ \ 
\includegraphics[scale=0.3]{images/CreativeCommonsSharealike}
\end{center}

\ \\

%\setbeamertemplate{itemize items}[$+$]
\begin{itemize}
\item[\textbullet] This presentation was created and narrated by Lane Schwartz. \\ \ \\
\item[\textbullet] You are free to reproduce and adapt this work under the terms of the Creative Commons Attribution-ShareAlike 4.0 International License.
\end{itemize}

\end{frame}


%\begin{frame}
\frametitle{Linear unit}

\begin{equation*}
y = b + \sum_i x_i w_i
\end{equation*}


\begin{tikzpicture}
\begin{axis}[ylabel=output,xlabel=weighted input,scale=0.3,grid=major,ytick={0},xtick={0},xmin=-5,xmax=5,ymin=-5,ymax=5]
  \addplot[blue, ultra thick][domain=-10:10] (x,x);
\end{axis}
\end{tikzpicture}

\end{frame}


%\frame {

\begin{center}
$f(x_1,x_2) = \sin(x_1) + x_1*x_2$ \ \ \ \ \ \ \ \ \ \ 
$\bar{x_i} = \frac{\partial f_j}{\partial x_i}$
\pause

\begin{tikzpicture}[>=latex, scale=0.75, transform shape]
\tikzstyle{every node}=[font=\large]
\tikzstyle{every path}=[line width=1pt]

\node[state,rectangle] (fxx) {$f(x_1,x_2)$};
\node[state] (plus) [below=1.5cm of fxx] {$f_5$};
\node[state] (sin) [below left=1.5cm and 0.5cm of plus] {$f_4$};
\node[state] (mul) [below right=1.5cm and 0.5cm of plus] {$f_3$};
\node[state] (x1) [below left=1.5cm and 0.5 cm of sin] {$f_1$};
\node[state] (x2) [below right=1.5cm and 0.5 cm of mul] {$f_2$};

\onslide<+>{
\path[-]
(plus) edge node {} (fxx)
(sin) edge node {} (plus)
(mul) edge node {} (plus)
(x1) edge node {} (sin)
(x1) edge node {} (mul)
(x2) edge node {} (mul);
}

\onslide<+>{
\path[->,blue]
(plus) edge node {} (fxx)
(sin) edge node {} (plus)
(mul) edge node {} (plus)
(x1) edge node {} (sin)
(x1) edge node {} (mul)
(x2) edge node {} (mul);
}

\onslide<.->{
\node[state, draw=none,fill=none,blue] (w1) [left=0.1cm of x1] {$y_1 = f_1(x_1) = x_1$};
\node[state, draw=none,fill=none,blue] (w2) [right=0.1cm of x2] {$y_2 = f_2(x_2) = x_2$};
\node[state, draw=none,fill=none,blue] (w3) [right=0.1cm of mul] {$y_3 = f_3(y_1, y_2) = y_1 \cdot y_2$};
\node[state, draw=none,fill=none,blue] (w4) [left=0.1cm of sin] {$y_4 = f_4(y_1) = \sin(y_1)$};
\node[state, draw=none,fill=none,blue] (w5) [left=0.1cm of plus] {$y_5 = f_5(y_3, y_4) = y_3 + y_4$};
\node[state, draw=none,fill=none,blue] (w6) [right=0.1cm of fxx] {$y_6 = f_6(y_5) = y_5$};
}

\onslide<+>{
\path[<-,red]
(plus) edge node [right] {$\bar{f_6} = \bar{y}_5 = \frac{\partial f_6}{\partial y_5} = 1$} (fxx)
(sin) edge node [left, pos=.6] {$\bar{y}_4 = \frac{\partial f_6}{\partial y_4} = \frac{\partial f_6}{\partial y_5} \frac{\partial y_5}{\partial y_4} = \bar{y}_5 \frac{\partial f_5}{\partial y_4} = \bar{y}_5 \cdot 1 = \bar{y}_5$} (plus)
(mul) edge node [right, pos=.6] {$\bar{y}_3 = \frac{\partial f_6}{\partial y_3} = \frac{\partial f_6}{\partial y_5} \frac{\partial y_5}{\partial y_3} = \bar{y}_5 \frac{\partial f_5}{\partial y_3} = \bar{y}_5 \cdot 1 = \bar{y}_5$} (plus)
(x1) edge node [left,pos=.6] { $\frac{\partial f_6}{\partial y_1} = \bar{y}_4 \frac{\partial y_4}{\partial y_1} = \bar{y}_4\cdot\cos(y_1)$ } (sin)
(x1) edge node [right, pos=.3] {$\bar{y}_{1^b}= \bar{y}_3\cdot y_2$} (mul)
(x2) edge node [right, pos=.6] {$\bar{y}_2 = \bar{y}_3 \frac{\partial y_3}{\partial y_2} = \bar{y}_3\cdot y_1$} (mul);

}

\end{tikzpicture}
\end{center}
}




\frame {

\begin{center}
$f(x_1,x_2) = \sin(x_1) + x_1*x_2$

\ \\

\pause

\begin{tikzpicture}[>=latex, scale=0.75, transform shape]
\tikzstyle{every node}=[font=\large]
\tikzstyle{every path}=[line width=1pt]

\node[state,rectangle] (fxx) {$f(x_1,x_2)$};
\node[state] (plus) [below=1.5cm of fxx] {$+$};
\node[state] (sin) [below left=1.5cm and 0.5cm of plus] {$\sin$};
\node[state] (mul) [below right=1.5cm and 0.5cm of plus] {$*$};
\node[state,rectangle] (x1) [below left=1.5cm and 0.5 cm of sin] {$x_1$};
\node[state,rectangle] (x2) [below right=1.5cm and 0.5 cm of mul] {$x_2$};

\onslide<+>{
\path[-]
(plus) edge node {} (fxx)
(sin) edge node {} (plus)
(mul) edge node {} (plus)
(x1) edge node {} (sin)
(x1) edge node {} (mul)
(x2) edge node {} (mul);
}

\onslide<+>{
\path[->,blue]
(plus) edge node {} (fxx)
(sin) edge node {} (plus)
(mul) edge node {} (plus)
(x1) edge node {} (sin)
(x1) edge node {} (mul)
(x2) edge node {} (mul);
}

\onslide<.->{
\node[state, draw=none,fill=none,blue] (w1) [left=0.1cm of x1] {$y_1 = x_1$};
\node[state, draw=none,fill=none,blue] (w2) [right=0.1cm of x2] {$y_2 = x_2$};
\node[state, draw=none,fill=none,blue] (w3) [right=0.1cm of mul] {$y_3 = y_1 \cdot y_2$};
\node[state, draw=none,fill=none,blue] (w4) [left=0.1cm of sin] {$y_4 = \sin(y_1)$};
\node[state, draw=none,fill=none,blue] (w4) [left=0.1cm of plus] {$y_5 = y_3 + y_4$};
}

\onslide<+>{
\path[<-,red]
(plus) edge node [right] {$\bar{f} = \bar{y}_5 = \frac{\partial f}{\partial y_5} = 1$} (fxx)
(sin) edge node [left, pos=.6] {$\bar{y}_4 = \bar{y}_5 \frac{\partial y_5}{\partial y_4}$} (plus)
(mul) edge node [right, pos=.6] {$\bar{y}_3 = \bar{y}_5 \frac{\partial y_5}{\partial y_3}$} (plus)
(x1) edge node [left,pos=.6] { $\bar{y}_1^a = \bar{y}_4 \frac{\partial y_4}{\partial y_1} = \bar{y}_4\cdot\cos(y_1)$ } (sin)
(x1) edge node [right, pos=.3] {$\bar{y}_1^b= \bar{y}_3\cdot y_2$} (mul)
(x2) edge node [right, pos=.6] {$\bar{y}_2 = \bar{y}_3 \frac{\partial y_3}{\partial y_2} = \bar{y}_3\cdot y_1$} (mul);

}


\end{tikzpicture}
\end{center}
}



\end{document}
